\chapter{Canonical Data Schema}\label{ch:schema}

\section{Overview}

Every artefact under review is represented by a JSON record validated
against a Draft-07 JSON Schema. Six schemas are defined; collectively
they constitute the data contract consumed by the AI pipeline
(\cref{ch:ai}) and by the controls engine (\cref{ch:controls-engine}).

\begin{table}[htbp]
\caption{Schemas under \texttt{structured/schema/}.}\label{tab:schemas}
\centering
\small
\begin{tabularx}{\linewidth}{@{}lXr@{}}
\toprule
\textbf{Schema} & \textbf{Validates} & \textbf{Records} \\ \midrule
\path{control.schema.json}         & EUC and process controls            & 3 \\
\path{decision-point.schema.json}  & Gateway logic records                & 5 \\
\path{workflow.schema.json}        & State machine for a TB workflow      & 1 \\
\path{requirement.schema.json}     & Business-case requirement            & 1 \\
\path{kpi.schema.json}             & KPI definitions (collection file)    & 6 \\
\path{glossary.schema.json}        & Domain glossary                      & 1 (28 terms) \\
\path{data-product.schema.json}    & ODPS 4.1 data-product manifest       & 1 \\
\bottomrule
\end{tabularx}
\end{table}

\section{Control record}\label{sec:control-schema}

A \texttt{control} binds a governance reference (the EUC identifier) to
one or more workflow states. It declares type, frequency, owner, and
the evidence location examined at audit.

\begin{table}[htbp]
\caption{\texttt{control.schema.json} fields.}\label{tab:control-fields}
\centering
\small
\begin{tabularx}{\linewidth}{@{}lllX@{}}
\toprule
\textbf{Field} & \textbf{Type} & \textbf{Req} & \textbf{Notes} \\ \midrule
controlId        & string  & Y & \texttt{EUC-\textit{nnn}} \\
eucReference     & string  & N & \texttt{EX/nnn/nnnnnnn/nnnn/nnn} pattern \\
controlObjective & string  & Y & Plain-language objective \\
controlType      & enum    & Y & \textsc{preventive}, \textsc{detective}, \textsc{corrective} \\
frequency        & enum    & Y & \textsc{daily}..\textsc{annually} \\
timing           & string  & N & \texttt{M$\pm$n} phase \\
owner            & string  & Y & GFS role or named individual \\
evidenceLocation & string  & N & Where audit evidence lives \\
relatedStates    & array   & Y & Workflow \texttt{stateId}s ($\ge$1) \\
workflowRef      & string  & N & Workflow \texttt{workflowId} \\
source\_document & object  & N & Path, optional page and sha256 \\
\bottomrule
\end{tabularx}
\end{table}

\begin{figure}[htbp]
\begin{lstlisting}[language=jsonx,caption={\texttt{EUC-002} variance-file
control record (abbreviated).},label=lst:euc-002]
{
  "controlId": "EUC-002",
  "eucReference": "EX/200/1579702/0001/000",
  "controlObjective": "TB Variance Analysis file integrity and calculation accuracy",
  "controlType": "detective",
  "frequency": "monthly",
  "timing": "M-2 to M+5",
  "owner": "GFS Analyst",
  "evidenceLocation": "Shared Drive",
  "relatedStates": ["UPDATE_TB_DETAILS", "CALCULATE_VARIANCE", "FILTER_VARIANCES"],
  "description": "Ensures variance calculations in the Excel working file are accurate."
}
\end{lstlisting}
\end{figure}

\section{Decision-point record}\label{sec:decision-point-schema}

A decision point encodes an exclusive gateway: an input metric, an
operator, a threshold (where applicable), and an enumerated list of
outcomes. Each outcome pairs a human-readable condition with a
machine-actionable \texttt{action} verb. The \texttt{aiAugmentation}
block names the substitution target for the manual judgement.

\begin{figure}[htbp]
\begin{lstlisting}[language=jsonx,caption={\texttt{DP-001} material
variance gateway.},label=lst:dp-001]
{
  "id": "DP-001",
  "name": "Material Variance Check",
  "stateRef": "CHECK_MATERIAL_VARIANCES",
  "inputMetric": "variance_amount",
  "operator": "exceeds_threshold",
  "threshold": "materiality_limit",
  "outcomes": [
    { "condition": "Variances Are Material",  "action": "seek_commentaries" },
    { "condition": "Variances Not Material",  "action": "proceed_to_unexplained_check" }
  ],
  "aiAugmentation": {
    "opportunity": "Auto-classify variance materiality using ML threshold tuning.",
    "dataNeeded": ["variance_amount", "account_category", "historical_thresholds"]
  }
}
\end{lstlisting}
\end{figure}

\section{Workflow record}\label{sec:workflow-schema}

The workflow record declares a state machine. Valid \texttt{type}
values are \textsc{start\_event}, \textsc{end\_event}, \textsc{task},
\textsc{gateway\_exclusive}, \textsc{gateway\_parallel}, and
\textsc{intermediate\_event}. Each state carries optional
\texttt{systems}, \texttt{inputs}, \texttt{outputs}, a
\texttt{participant}, and a list of \texttt{transitions}. A transition
in a gateway state additionally carries a \texttt{condition} string
that matches one of the outcome labels in the associated decision
point.

\begin{table}[htbp]
\caption{Workflow state-machine header fields.}\label{tab:workflow-fields}
\centering
\small
\begin{tabularx}{\linewidth}{@{}lllX@{}}
\toprule
\textbf{Field} & \textbf{Type} & \textbf{Req} & \textbf{Notes} \\ \midrule
workflowId    & string & Y & Canonical slug (\texttt{tb-review-glo}) \\
version       & string & Y & SemVer \\
title         & string & Y & Long name \\
processOwner  & string & N & Named owner \\
timeline      & array  & N & Phase / timing / description \\
states        & array  & Y & $\ge 2$ states, including one \textsc{start\_event} and one \textsc{end\_event} \\
dataObjects   & array  & N & Workflow-scoped artefacts \\
participants  & array  & N & Roles and groups \\
\bottomrule
\end{tabularx}
\end{table}

\section{Requirement record}\label{sec:requirement-schema}

The requirement record captures the business case. It is scoped by a
status enum (\textsc{draft}, \textsc{review}, \textsc{approved},
\textsc{rejected}, \textsc{superseded}), contains executive summary,
scope, stakeholders, process steps (each with an \texttt{aiPriority}
enum from \textsc{low}..\textsc{critical}), decision points, KPI
references, control references, risk assessment, and a phased
implementation plan. The single record populated here is
\texttt{TB-REQ-001} in state \textsc{approved}.

\section{KPI record}\label{sec:kpi-schema}

KPIs are stored as a collection file at
\corpus\texttt{structured/kpis/tb-kpis.json} containing one array
entry per KPI (\cref{tab:kpis}). Each entry is validated against
\path{kpi.schema.json}; the wrapper is checked by
\path{structured/validate.py}.

\begin{table}[htbp]
\caption{The six KPIs defined for the TB process.}\label{tab:kpis}
\centering
\small
\begin{tabularx}{\linewidth}{@{}llXl@{}}
\toprule
\textbf{KPI} & \textbf{Type} & \textbf{Description} & \textbf{Target} \\ \midrule
TB-KPI-001 & threshold  & Variance materiality threshold       & country-specific \\
TB-KPI-002 & duration   & TB review cycle time                 & reduce via daily controls \\
TB-KPI-003 & percentage & Commentary coverage                  & 100 \\
TB-KPI-004 & effort     & FTE effort                           & 15--20 FTEs \\
TB-KPI-005 & percentage & Unexplained variance rate            & 0 \\
TB-KPI-006 & percentage & EUC control compliance               & 100 \\
\bottomrule
\end{tabularx}
\end{table}

\section{Glossary record}

One glossary record contains 28 terms drawn verbatim from the
business case and the DOI. Category membership
(\textsc{balance\_sheet}, \textsc{income\_statement}, \textsc{general},
\textsc{regulatory}) is used by the GlossaryService pair-type alias
expansion at query time.

\section{Data-product record}

The data-product record is an ODPS 4.1-conformant manifest. It names
the routing policy (\texttt{vllm-only}), the data security class
(\textsc{confidential}), the owner, the per-country view for HK, the
system prompt, and the RAG source wiring. Consumers read this record
to know whether a given user request is permitted to contact external
LLM providers --- the answer for all TB artefacts is \emph{no}.

\section{Enumerations}

All enumerated values used in any schema or record are listed
authoritatively in \corpus\texttt{structured/enums.yaml}.
\Cref{tab:enums} summarises the most frequently referenced groups.

\begin{table}[htbp]
\caption{Frequently referenced enumerations.}\label{tab:enums}
\centering
\small
\begin{tabularx}{\linewidth}{@{}lX@{}}
\toprule
\textbf{Enum} & \textbf{Values} \\ \midrule
control\_type        & preventive, detective, corrective \\
workflow\_state\_type & start\_event, end\_event, task, gateway\_exclusive, gateway\_parallel, intermediate\_event \\
decision\_operator   & equals, not\_equals, exceeds\_threshold, within\_threshold, in\_set, boolean \\
ai\_priority         & low, medium, high, critical \\
risk\_severity       & L, M, H \\
kpi\_measure\_type    & threshold, duration, percentage, effort, count, ratio \\
data\_security\_class & public, internal, confidential, restricted \\
\bottomrule
\end{tabularx}
\end{table}

\section{Validation}

A local registry-based validator is provided at
\corpus\texttt{structured/validate.py}. It walks the record
directories, maps each to its schema via
\path{RECORD_DIRS}, and reports per-file results. The script is
offline: it instantiates a \texttt{referencing.Registry} loaded from
\corpus\texttt{structured/schema/} so no network resolution is
attempted.

\begin{lstlisting}[caption={Running the validator.},label=lst:validate]
$ python3 docs/tb/structured/validate.py
[ok]   controls/euc-001-tb-data-extraction.json
[ok]   controls/euc-002-variance-analysis-file.json
[ok]   controls/euc-003-internal-review.json
[ok]   decision-points/dp-001-material-variance.json
...
[ok]   corpus.json (17 items, all paths resolved)
13/13 records validated, corpus OK
\end{lstlisting}

\section{Variance record}\label{sec:variance-schema}

A variance record is produced by the Stage 2 anomaly detection pipeline
(\cref{ch:ai}). It captures the statistical deviation of an account
balance from its historical trend and records the outcome of initial
materiality and anomaly checks.

\begin{table}[htbp]
\caption{\texttt{variance-record.schema.json} fields.}\label{tab:variance-fields}
\centering
\small
\begin{tabularx}{\linewidth}{@{}lllX@{}}
\toprule
\textbf{Field} & \textbf{Type} & \textbf{Req} & \textbf{Notes} \\ \midrule
\texttt{variance\_id}       & string  & Y & Unique record identifier \\
\texttt{account\_code}      & string  & Y & GL account code \\
\texttt{period}             & string  & Y & YYYY-MM format \\
\texttt{current\_balance}   & number  & N & Current period balance \\
\texttt{prior\_balance}     & number  & N & Prior period balance \\
\texttt{variance\_amount}   & number  & Y & Absolute variance ($|C - P|$) \\
\texttt{variance\_percent}  & number  & N & Percentage variance \\
\texttt{flags}              & array   & Y & Anomaly flags (\textsc{spike}, \textsc{flip}, \dots) \\
\texttt{anomaly\_detection}  & object  & N & Statistical check details (3$\sigma$, MAD) \\
\texttt{materiality\_assessment} & object  & N & DP-001 assessment details \\
\texttt{investigation\_status}   & enum    & N & \textsc{pending}, \textsc{in\_progress}, \textsc{completed} \\
\texttt{is\_unexplained}     & boolean & N & True if remains unexplained at M+4 \\
\bottomrule
\end{tabularx}
\end{table}

\section{Unified Schema Registry}

As of 18 April 2026, all schemas in this specification are published in the
unified schema registry at \path{docs/schema/}. Migration status is
\emph{complete} (see \path{docs/schema/registry.json} field
\texttt{migration.status}). This consolidation addresses cross-document schema
divergence identified in the specification review.

\begin{description}
  \item[Registry index.] \path{docs/schema/registry.json} --- Master index of
        all schemas across domains.
  \item[Common types.] \path{docs/schema/common/base-types.schema.json} ---
        Shared definitions (\texttt{Metadata}, \texttt{Currency},
        \texttt{DateRange}) referenced via JSON Schema \texttt{\$ref}.
  \item[Consolidated enums.] \path{docs/schema/common/enums.yaml} --- Single
        source of truth for enumeration values including \texttt{account\_category},
        \texttt{variance\_flag}, \texttt{decision\_point}, and
        \texttt{workflow\_status}.
  \item[TB schemas.] \path{docs/schema/tb/} --- Domain-specific schemas:
        \texttt{tb-extract.schema.json}, \texttt{pl-extract.schema.json},
        \texttt{variance-record.schema.json}, \texttt{commentary-draft.schema.json},
        and \texttt{decision-point.schema.json}. All marked \texttt{status: stable}.
\end{description}

\noindent
The legacy YAML schemas at \path{docs/tb/machine-readable/training-data/hkg-tb-schema.yaml}
and \path{docs/tb/machine-readable/training-data/hkg-pl-schema.yaml} have been
converted to JSON Schema Draft-07 in the unified registry; the registry
entries record the \texttt{legacy\_path} for audit traceability. The legacy
files are retained read-only during a one-release deprecation window and
will be removed in v1.1. Validation:

\begin{lstlisting}[language=bash]
python3 docs/schema/validate.py --registry-check
python3 docs/schema/validate.py --domain tb \
    --schema tb-extract.schema.json --file tb-hkg-2026-03.json
\end{lstlisting}

\noindent
CI runs \texttt{--registry-check} on every pull request that touches
\path{docs/schema/} or any TB pipeline module.
